% Source printed pp. 210-214 / PDF pp. 218-222.
\section{光滑曲线回顾}\label{src:6.1}
\begin{epigraph}
每个人都知道曲线是什么, 直到他学了足够多的数学, 开始为无数可能的例外而困惑.
\epigraphauthor{Felix Klein}
\end{epigraph}
我们回顾光滑曲线的基本性质, 下一节将把它们推广到结点曲线.
\subsection{曲线}\label{src:6.1.1-subsection}
\begin{definition}\label{src:6.1.1}
域 $k$ 上的\newterm{曲线}{curve} 是 $k$ 上的一维有限型概形 $C$.
\end{definition}
固有曲线是射影的~\cite[命题 I.6.7]{Har77}; 而且, 每个一维固有代数空间事实上都是射影概形 (定理~\bookref{5.5.32}). 若 $C$ 是域 $k$ 上的固有曲线, 定义 $C$ 的\newterm{算术亏格}{arithmetic genus}, 或简称\textit{亏格}, 为
\[
g(C)=1-\chi(C,\mathcal O_C).
\]
若 $C$ 几何连通且约化, 它等于 $\mathrm h^1(C,\allowbreak \mathcal O_C)$. 约化固有曲线 $C$ 的\newterm{几何亏格}{geometric genus} 是其正规化 $\widetilde C$ 的 (算术) 亏格.

对代数闭域 $k$ 上的约化, 连通射影曲线 $C$, $C$ 上甚丰沛线丛 $L$ 的次数定义为 $L$ 的任意截面的零点个数 (计重数). 换言之, 若 $C\hookrightarrow\mathbb P^n$ 是 $L$ 定义的射影嵌入, 则 $\deg L=\dim_k\Gamma(C\cap H,\allowbreak \mathcal O_{C\cap H})$, 其中 $H$ 是任意超平面, $C\cap H$ 是概形意义的交. $C$ 上的任意线丛都能写成两个甚丰沛线丛的差: 若 $M$ 在 $C$ 上甚丰沛, 则对 $n\gg0$, $M':=L\otimes M^n$ 甚丰沛, 且 $L\cong M'\otimes(M^{\otimes n})^\vee$. 由此可定义 $\deg L=\deg M'-n\deg M$. 注意次数对张量积可加; 若 $C=\bigcup_i C_i$ 是不可约分解, 则 $\deg L=\sum_i\deg L|_{C_i}$.
\begin{theorem}[Riemann--Roch]\label{src:6.1.2}
设 $C$ 是代数闭域 $k$ 上亏格为 $g$ 的约化, 连通射影曲线. 若 $L$ 是 $C$ 上的线丛, 则
\[
\chi(C,L)=\deg L+1-g.
\]
\end{theorem}
\begin{proof}
可写 $L=\mathcal O_C(D)$, 其中除子 $D$ 支撑在光滑轨迹上. 由亏格的定义, Riemann--Roch 对 $\mathcal O_C$ 成立; 因而通过增减点, 只需证明: 对光滑点 $p\in C(k)$, Riemann--Roch 对 $\mathcal O_C(D)$ 成立当且仅当对 $\mathcal O_C(D+p)$ 成立. 考虑短正合列
\[
0\to\mathcal O_C(D)\to\mathcal O_C(D+p)\to\kappa(p)\to0
\]
及恒等式 $\chi(C,\allowbreak \mathcal O_C(D+p))=\chi(C,\allowbreak \mathcal O_C(D))+1$, 即得结论. 另见~\cite[定理 IV.1.3, 习题 IV.1.9]{Har77} 和~\cite[习题 18.4.B 和 S]{Vak17}.
\end{proof}
\subsection{光滑曲线}\label{src:6.1.2-subsection}
若 $C$ 是光滑曲线, 则微分层 $\Omega_C$ 是线丛. Serre 对偶断言 $\Omega_C$ 是 $C$ 上的对偶化层, 是研究曲线时不可或缺的深刻结果.
\begin{theorem}[光滑曲线的 Serre 对偶]\label{src:6.1.3}
若 $C$ 是域 $k$ 上的光滑射影曲线, 则 $\Omega_C$ 是对偶化层, 即存在线性映射 $\operatorname{tr}:\mathrm H^1(C,\allowbreak \Omega_C)\to\allowbreak k$, 使得对每个凝聚层 $F$, 自然配对
\[
\operatorname{Hom}_{\mathcal O_C}(F,\Omega_C)\times\mathrm H^1(C,F)\to\mathrm H^1(C,\Omega_C)\xrightarrow{\operatorname{tr}}k
\]
是完美配对.
\end{theorem}
\begin{proof}
见~\cite[定理 II.8.2]{Ser88} 和~\cite[推论 III.7.12]{Har77}.
\end{proof}
\begin{remark}\label{src:6.1.4}
配对完美是指 $\operatorname{Hom}_{\mathcal O_C}(F,\allowbreak \Omega_C)$ 等同于对偶 $\mathrm H^1(C,\allowbreak F)^\vee$. 若 $F$ 是向量丛, 则 $\operatorname{Hom}_{\mathcal O_C}(F,\allowbreak \Omega_C)\cong\mathrm H^0(C,\allowbreak F^\vee\otimes\Omega_C)$, Serre 对偶给出同构
\[
\mathrm H^0(C,F^\vee\otimes\Omega_C)\cong\mathrm H^1(C,F)^\vee.
\]
取 $F=\Omega_C$, 可见 $\mathrm H^1(C,\allowbreak \Omega_C)\cong\mathrm H^0(C,\allowbreak \mathcal O_C)^\vee$; 若 $C$ 几何连通且约化, 则 $\operatorname{tr}:\mathrm H^1(C,\allowbreak \Omega_C)\to\allowbreak k$ 是同构.
\end{remark}
Serre 对偶给出更强的 Riemann--Roch 形式.
\begin{theorem}[Riemann--Roch II]\label{src:6.1.5}
设 $C$ 是代数闭域上亏格为 $g$ 的光滑, 连通射影曲线. 若 $L$ 是 $C$ 上的线丛, 则
\[
\mathrm h^0(C,L)-\mathrm h^0(C,\Omega_C\otimes L^\vee)=\deg L+1-g.\qquad\square
\]
\end{theorem}
\begin{remark}\label{src:6.1.6}
这常写成除子形式 $\mathrm h^0(C,\allowbreak L)-\mathrm h^0(C,\allowbreak K-L)=\deg L+1-g$, 其中 $K$ 是典范除子, 即 $\Omega_C=\mathcal O_C(K)$.
\end{remark}
Riemann--Hurwitz 定理 (\bookref{6.7.4}) 也在光滑曲线的研究中起着重要作用, 它描述微分层在有限覆盖下的行为. 我们将其推迟到 \S\bookref{6.7.1-subsection} 讨论分歧覆盖时处理.
\begin{exercise}[容易]\label{src:6.1.7}
对任意 $g\geq0$, 证明存在亏格为 $g$ 的光滑, 连通射影曲线.
\end{exercise}
\paragraph{自同构.}
若 $C$ 是域 $k$ 上的光滑曲线, $p_1,\allowbreak \ldots,\allowbreak p_n\in C(k)$, 则用 $\operatorname{Aut}(C,\allowbreak p_i)$ 表示 $(C,\allowbreak p_i)$ 在 $k$ 上的自同构的抽象群, 它由满足对所有 $i$ 都有 $\alpha(p_i)=p_i$ 的自同构 $\alpha:C\xrightarrow{\sim}C$ 组成.
\begin{proposition}\label{src:6.1.8}
设 $C$ 是代数闭域 $k$ 上亏格为 $g$ 的光滑, 连通射影曲线, $p_1,\allowbreak \ldots,\allowbreak p_n\in C(k)$ 是互异点. 若 $2g-2+n>0$, 即 $g\geq2$, 或 $g=1$ 且 $n\geq1$, 或 $g=0$ 且 $n\geq3$, 则自同构群 $\operatorname{Aut}(C,\allowbreak p_i)$ 有限. 此外, 若 $g\geq2$ 且 $\operatorname{char}(k)=0$, 则 $|\operatorname{Aut}(C)|\leq84(g-1)$.
\end{proposition}
\begin{proof}
这是 Hurwitz 的经典结果~\cite{Hur92}. 关于有限性, 另见~\cite[习题 IV.5.2]{Har77}, \cite[习题 I.F.4]{ACGH85} 和~\cite[定理 4.18]{Mir95}; 关于明确的界, 见~\cite[习题 IV.2.5]{Har77}, \cite[习题 I.F.10]{ACGH85} 和~\cite[定理 3.9]{Mir95}.
\end{proof}
\begin{exercise}\label{src:6.1.9}
设 $C$ 是代数闭域 $k$ 上亏格 $g\geq2$ 的光滑, 连通射影曲线. 证明 $C$ 不存在固定超过 $2g+2$ 个点的非平凡自同构.
\end{exercise}
每条超椭圆曲线 (因而每条亏格为 $2$ 的曲线) 都有非平凡的超椭圆对合, 但以下习题将在后文用于说明更高亏格的一般曲线没有非平凡自同构.
\begin{exercise}\label{src:6.1.10}
证明若 $g\geq3$ 且 $k$ 是代数闭域, 则存在 $k$ 上自同构群平凡的光滑, 连通射影曲线 $C$.
\end{exercise}
关于光滑曲线的更多背景, 推荐~\cite[\S IV]{Har77}, \cite[\S 20]{Vak17}, \cite[\S 7]{Liu02}, \cite{ACGH85} 和~\cite{Mir95}.
\subsection{光滑曲线上除子的正性}\label{src:6.1.3-subsection}
Riemann--Roch 的以下推论给出判断线丛是否无基点 (等价于整体生成), 丰沛或甚丰沛的实用判据.
\begin{corollary}\label{src:6.1.11}
设 $C$ 是代数闭域 $k$ 上的光滑, 连通射影曲线, $L$ 是 $C$ 上的线丛.
\begin{enumerate}[(1)]
\item 若 $\deg L<0$, 则 $\mathrm h^0(C,\allowbreak L)=0$;
\item 若 $\deg L>0$, 则 $L$ 丰沛;
\item 若 $\deg L\geq2g$, 则 $L$ 无基点; 且
\item 若 $\deg L\geq2g+1$, 则 $L$ 甚丰沛.
\end{enumerate}
\end{corollary}
\begin{proof}
见~\cite[推论 IV.3.2]{Har77}.
\end{proof}
\begin{remark}\label{src:6.1.12}
若 $g>1$, 可用 Riemann--Roch 和 Serre 对偶计算得: (a) $\mathrm h^0(C,\allowbreak \Omega_C)=\mathrm h^1(C,\allowbreak \mathcal O_C)=g$, (b) $\mathrm h^1(C,\allowbreak \Omega_C)=\mathrm h^0(C,\allowbreak \mathcal O_C)=1$, 且 (c) $\Omega_C$ 的次数为 $2g-2$, 因而在 $C$ 上丰沛. 类似地, 若 $k>1$, 则: (a) $\mathrm h^0(C,\allowbreak \Omega_C^{\otimes k})=(2k-1)(g-1)$, (b) $\mathrm h^1(C,\allowbreak \Omega_C^{\otimes k})=0$, 且 (c) $\Omega_C^{\otimes k}$ 的次数为 $2k(g-1)$, 若 $k\geq3$ 则甚丰沛. 注意 $\Omega_C$ 不甚丰沛当且仅当 $C$ 是超椭圆曲线. 另一方面, 若 $g=1$, 则 $\Omega_C\cong\mathcal O_C$; 若 $g=0$, 则 $C=\mathbb P^1$, $\Omega_C=\mathcal O(-2)$.
\end{remark}
\subsection{有理曲线的分类}\label{src:6.1.4-subsection}
在代数闭域 $k$ 上, 每条亏格为 $0$ 的连通, 光滑固有曲线都同构于 $\mathbb P_k^1$. 在任意域上, 分类略为复杂.
\begin{exercise}[射影直线的分类]\label{src:6.1.13}
证明对域 $k$ 上的固有整曲线 $C$, 以下条件等价:
\begin{enumerate}[(a)]
\item $C\cong\mathbb P_k^1$;
\item $C$ 在 $k$ 上光滑且几何不可约, 亏格为 $0$, 且 $C(k)\neq\varnothing$;
\item $C$ 是 Gorenstein 的且在 $k$ 上几何整, 亏格为 $0$, 且具有奇数次线丛;
\item $\mathrm H^0(C,\allowbreak \mathcal O_C)=k$, $C$ 的亏格为 $0$, 且具有次数为 $1$ 的线丛;
\item $\mathrm H^1(C,\allowbreak \mathcal O_C)=0$, 且 $C$ 具有次数为 $1$ 的线丛.
\end{enumerate}
另见~\cite[Tag 0C6U]{SP}.
\end{exercise}
还有奇异有理曲线的分类, 它有助于理解结点有理曲线及预稳定曲线中有理尾和有理桥的分类.
\begin{exercise}[奇异有理曲线的分类]\label{src:6.1.14}
设 $C$ 是域 $k$ 上亏格为 $0$ 的奇异固有整曲线, 满足 $\mathrm H^0(C,\allowbreak \mathcal O_C)=k$, 设 $\pi:C'\to\allowbreak C$ 是正规化. 证明:
\begin{enumerate}[(a)]
\item $C$ 有唯一奇异点 $p$, 它是 $k$-有理点;
\item $\pi:C'\to\allowbreak C$ 在 $C\setminus p$ 上是同构, 且 $\pi^{-1}(p)=\{p'\}$, 其中 $p'\in C(k')$, $k'/\mathrm{base}$ 是有限域扩张;
\item $C'\cong\mathbb P_{k'}^1$; 且
\item $C$ 等同于 Ferrand 推出 (\bookref{B.4.1})
\[
\begin{tikzcd}
\Spec k'\arrow[r,"p'"]\arrow[d]&C'\arrow[d,"\pi"]\\
\Spec k\arrow[r,"p"']\arrow[ur,phantom,"\lrcorner",very near start]&C.
\end{tikzcd}
\]
\end{enumerate}
另见~\cite[Tag 0DJB]{SP}.
\end{exercise}
\subsection{光滑曲线族}\label{src:6.1.5-subsection}
\begin{definition}\label{src:6.1.15}
概形 $S$ 上的\newterm{光滑曲线族}{family of smooth curves} (亏格为 $g$) 是概形的光滑固有态射 $C\to\allowbreak S$, 使得每个几何纤维都是连通曲线 (亏格为 $g$).
\end{definition}
回顾相对微分层 $\Omega_{C/S}$ 是 $C$ 上的线丛, 对 $s\in S$, 限制 $\Omega_{C/S}|_{C_s}$ 等同于 $\Omega_{C_s}$. 更一般地, 对每个概形态射 $T\to\allowbreak S$, $\Omega_{C/S}$ 在 $C_T:=C\times_S T$ 上的拉回典范同构于 $\Omega_{C_T/T}$. 推广推论~\bookref{6.1.11}, 对 $k\geq3$, $\Omega_{C/S}^{\otimes k}$ 相对甚丰沛, 且其推前是 $S$ 上的向量丛.
\begin{proposition}[光滑曲线族的性质]\label{src:6.1.16}
设 $\pi:C\to\allowbreak S$ 是亏格 $g\geq2$ 的光滑曲线族.
\begin{enumerate}[(1)]
\item $\pi_*\mathcal O_C=\mathcal O_S$;
\item 推前 $\pi_*(\Omega_{C/S}^{\otimes k})$ 是秩为
\[
r(k):=\begin{cases}g&k=1,\\(2k-1)(g-1)&k>1\end{cases}
\]
的向量丛, 其构造与基变换交换 (即对概形态射 $f:T\to\allowbreak S$, $f^*\pi_*(\Omega_{C/S}^{\otimes k})\cong\pi_{T,\allowbreak *}(\Omega_{C_T/T}^{\otimes k})$);
\item $R^1\pi_*\Omega_{C/S}^{\otimes k}$ 在 $k=1$ 时同构于 $\mathcal O_S$, 否则为零; 且
\item 对 $k\geq3$, $\Omega_{C/S}^{\otimes k}$ 相对甚丰沛. 特别地, $C\to\allowbreak S$ 射影.
\end{enumerate}
\end{proposition}
\begin{proof}
第 (1)--(3) 项由上同调与基变换 (\bookref{A.6.8}) 得到, 详见命题~\bookref{A.6.11}. 对 (4), 注意对每个点 $s\in S$, 纤维 $\Omega_{C/S}^{\otimes k}\otimes\kappa(s)=\Omega_{C_s}^{\otimes k}$ 由推论~\bookref{6.1.11} 甚丰沛, 因为 $\deg\Omega_{C_s}^{\otimes k}=k(2g-2)>0$. 由于 $\mathrm H^1(C_s,\allowbreak \Omega_{C_s}^{\otimes k})=0$, 可应用命题~\bookref{B.2.10}, 得到 $\Omega_{C/S}^{\otimes k}$ 相对甚丰沛.
\end{proof}
相对微分层 $\Omega_{C/S}$ 也是相对对偶化层, 即满足相对版本的 Serre 对偶; 见~\cite{Har66c} 或~\cite[\S 6.4]{Liu02}.
